% =========================================================================
% SEÇÃO 1: RESUMO
% =========================================================================
\section{Resumo}\label{sec:resumo}

Fluxos em digrafos constituem uma área fundamental da Otimização Combinatória, com amplas aplicações práticas em redes de transporte e alocação de recursos. Este projeto adota uma abordagem que combina fundamentos teóricos, implementação algorítmica e validação experimental para o estudo de problemas de fluxo máximo e fluxo de custo mínimo (veja~\cite{ahuja1993network}). Foram estudados e implementados algoritmos clássicos e modernos de fluxo máximo e de fluxo de custo mínimo, bem como técnicas de modelagem e reduções combinatórias, incluindo a aplicação prática em Segmentação Interativa de Imagens via corte mínimo. A validação prática foi realizada através de experimentos computacionais utilizando instâncias de \textit{benchmarks} reconhecidas da literatura (veja~\cite{Jensen2022}).
